\documentclass[11pt]{article}
\usepackage[T1]{fontenc}
\usepackage[utf8]{inputenc}
\usepackage{lmodern}
\usepackage{microtype}
\usepackage[margin=1in]{geometry}
\usepackage{amsmath,amssymb,amsthm,mathtools}
\usepackage{booktabs,longtable,array}
\usepackage{enumitem}
\usepackage{xcolor}
\usepackage{hyperref}
\usepackage{listings}
\usepackage{fancyhdr}
\usepackage{url}
\usepackage{seqsplit}
\usepackage{graphicx}

\hypersetup{
  colorlinks=true,
  linkcolor=blue!45!black,
  citecolor=green!35!black,
  urlcolor=blue!55!black,
  pdftitle={Causal-Branch Parameterization of Reidemeister-III Unlocking},
  pdfauthor={OpenAI, prepared for the ProveIt repository}
}
\setlist{itemsep=2pt,topsep=4pt}
\setlength{\parindent}{0pt}
\setlength{\parskip}{5pt}
\setlength{\headheight}{14pt}
\pagestyle{fancy}
\fancyhf{}
\lhead{ProveIt Unknot Recognition, Report 25}
\rhead{Causal-branch RIII search}
\cfoot{\thepage}

\newtheorem{theorem}{Theorem}[section]
\newtheorem{corollary}[theorem]{Corollary}
\newtheorem{proposition}[theorem]{Proposition}
\newtheorem{lemma}[theorem]{Lemma}
\newtheorem{claim}[theorem]{Claim}
\theoremstyle{definition}
\newtheorem{definition}[theorem]{Definition}
\newtheorem{example}[theorem]{Example}
\newtheorem{algorithmic}[theorem]{Algorithm}
\theoremstyle{remark}
\newtheorem{remark}[theorem]{Remark}
\newtheorem{warning}[theorem]{Warning}

\newcommand{\N}{\mathbb{N}}
\newcommand{\F}{\mathcal{F}}
\newcommand{\R}{\mathcal{R}}
\newcommand{\T}{\mathcal{T}}
\newcommand{\supp}{\operatorname{supp}}
\newcommand{\birth}{\operatorname{birth}}
\newcommand{\poly}{\operatorname{poly}}
\newcommand{\bigO}{O}
\newcommand{\Unk}{\mathsf{Unknot}}
\newcommand{\Unknown}{\mathsf{Unknown}}
\newcommand{\code}[1]{\texttt{#1}}
\newcommand{\hashcode}[1]{\texttt{\seqsplit{#1}}}

\lstset{
  basicstyle=\ttfamily\small,
  columns=fullflexible,
  breaklines=true,
  frame=single,
  numbers=left,
  numberstyle=\tiny,
  showstringspaces=false,
  tabsize=4
}

\title{\textbf{Causal-Branch Parameterization of Reidemeister-III Unlocking}\\
\large A restricted quasi-polynomial simplification theorem and an integration-ready search for ProveIt}
\author{OpenAI\\Prepared for Vladimir Reshetnikov's \texttt{ProveIt} repository}
\date{October 7, 2026\\Reviewed baseline: commit \texttt{23893144e3aafd074a23cfb633222ab7e03e6d2e}}

\begin{document}
\maketitle

\begin{abstract}
The maintained ProveIt unknot recognizer uses a fast exact Khovanov/Bar--Natan
fallback and a crossing-nonincreasing Reidemeister preprocessor.  Its current
Reidemeister-III (RIII) helper searches short sequences, but after the first
move it follows only the neighborhood of the most recent move.  The repository
records substantial empirical gains from this search, while explicitly leaving
general quasi-polynomial unknot recognition open.

This report isolates and solves the parameterized combinatorial problem behind
that helper.  For a bounded-local reversible rewrite system, we define the
\emph{birth number} $b$ of a first-unlocking trace to be the number of neutral
actions whose footprints are disjoint from all earlier footprints.  We prove:
(i) the action-footprint graph of every shortest first-unlocking trace, including
its terminal decreasing action, is connected; (ii) all birth actions can be
commuted to the beginning; and (iii) traces of length at most $k$ and birth
number at most $b$ can be enumerated in
\[
   N^{b+O(1)}(Ck)^k
\]
time and polynomial working memory, where $C$ depends only on the finite local
rule set.  In the fixed-dart model used by ProveIt, an RIII footprint has at most
24 darts and bounded local incidence, so the theorem applies with an absolute
constant.  Repeating the search between crossing-decreasing RI/RII phases gives
a sound unknot preprocessor.  It is quasi-polynomial on families with
$b,k=O(\log N)$, improves to $N^{O(\log\log N)}$ for $b=O(1)$ and
$k=O(\log N)$, and is polynomial for $b=O(1)$ and
$k=O(\log N/\log\log N)$.

A proof audit also finds and repairs a tempting false claim: shortest local
traces need not have connected chronological prefixes.  Two independent
preparatory branches may be joined only by a later move.  The birth parameter is
therefore necessary.  The bundle contains a dependency-free reference
implementation, strict counterexamples separating the current last-neighbor
search from accumulated-support search, a finite oracle comparison on 5,000
generated local systems, and an opt-in ProveIt adapter targeting the reviewed
\code{simplify.py} blob.  No theorem that arbitrary unknot diagrams have
logarithmic $b$ or $k$ is claimed, and the adapter has not been run through the
full repository regression suite in this environment.
\end{abstract}

\tableofcontents
\newpage

\section{Executive summary}
\label{sec:summary}

\subsection{What is new}

The principal new object is a causal parameter for local Reidemeister search.
At a stalled diagram, an RIII move is \emph{neutral}: it preserves crossing
number but can change which RI/RII reductions are available.  A sequence of
neutral moves is \emph{first-unlocking} when no RI/RII move is available before
the last neutral move, and at least one is available immediately afterwards.
Each RIII trial changes a bounded set of fixed darts.  We call that set its
\emph{footprint}.

A move is a \emph{birth} when its footprint is disjoint from the union of all
previous footprints.  The birth number $b$ counts independent causal branches.
All other moves touch the accumulated footprint.  The core normal-form theorem
says that every first-unlocking trace can be replaced, without increasing its
length or birth number, by one in which all births occur first.  This is not a
mere presentation convention: it converts a global search into $b$ global
choices followed by a bounded-local search.

For depth $k$ and birth budget $b$, the resulting deterministic search has the
form
\[
 \underbrace{N\cdots N}_{b\text{ global births}}
 \;\underbrace{(C b)(C(b+1))\cdots(Ck)}_{\text{connected extensions}},
\]
up to polynomial bookkeeping.  Thus its running time is bounded by
$N^{b+O(1)}(Ck)^k$.

\subsection{What this proves for unknot simplification}

Let $D$ be an $N$-crossing diagram.  Fix deterministic orders for RI/RII
exhaustion and for the complete birth-front tree, and consider the
crossing-nonincreasing scheme that repeatedly:
\begin{enumerate}
\item exhausts all available RI/RII reductions;
\item when stalled, searches for an RIII first-unlocking trace of depth at most
      $k$ and birth number at most $b$;
\item applies the trace and resumes RI/RII reduction.
\end{enumerate}
Every successful run is a checkable Reidemeister certificate that $D$ is an
unknot.  Completeness below is for the policy-defined class on which these fixed
choices reach the empty diagram; mere existence of some different successful
episode sequence would require outer backtracking and is not claimed.  There are at most $N$ successful unlock episodes, because every episode
is followed by a reduction removing at least one crossing.  Therefore the total
running time remains
\[
  N^{b+O(1)}(Ck)^k.
\]
This yields three clean regimes.

\begin{center}
\begin{tabular}{@{}lll@{}}
\toprule
Birth bound & Depth bound & Consequence \\
\midrule
$b=O(1)$ & $k=O(\log N/\log\log N)$ & polynomial time \\
$b=O(1)$ & $k=O(\log N)$ & $N^{O(\log\log N)}$ \\
$b=O(\log N)$ & $k=O(\log N)$ & $N^{O(\log N)}$ \\
\bottomrule
\end{tabular}
\end{center}

The last two are quasi-polynomial.  The first is a genuinely polynomial
restricted class.

\subsection{What this does not prove}

No argument in this report shows that every unknot diagram admits a
crossing-nonincreasing RIII/RI/RII simplification with logarithmic depth or birth
number.  General Reidemeister sequences may need crossing-increasing moves, and
Lackenby's polynomial upper bound on the total number of Reidemeister moves does
not imply the monotone local hypotheses used here.  The complete ProveIt
recognizer must retain its exact fallback.

The result should therefore be read as a new explicit-parameter and
quasi-polynomial \emph{island}: it identifies two explicit structural parameters
whose control is sufficient, supplies the complete search under those
parameters, and narrows the remaining geometric problem to proving or predicting
such control on important diagram families.

\subsection{Artifact map}

The accompanying archive contains:
\begin{itemize}
\item this article in LaTeX and PDF;
\item \code{prototype/causal\_search.py}, a generic executable model;
\item \code{tests/test\_causal\_search.py}, including two strict separation
      fixtures;
\item \code{experiments/exhaustive\_validation.py} and the resulting JSON from
      5,000 deterministic generated systems;
\item \code{experiments/exhaustive\_catalog.py}, covering a finite catalog of
      608,400 three-site instances without sampling;
\item \code{integration/clustered\_r3\_snippet.py}, an opt-in adapter for the
      private ProveIt \code{\_Darts} representation;
\item a review patch and explicit integration checklist;
\item a claims ledger distinguishing proved, tested, untested, and open items.
\end{itemize}

\section{Reviewed baseline, motivation, and scope}
\label{sec:baseline}

\subsection{The maintained ProveIt pipeline}

The reviewed source is the public ProveIt repository at commit
\hashcode{23893144e3aafd074a23cfb633222ab7e03e6d2e}.  At that commit, the project
combines cheap diagram reductions and obstructions with exact Khovanov/Bar--Natan
computation.  The commit message records 431 maintained tests and states
explicitly that general quasi-polynomial recognition remains unproved.  This
report accepts that boundary \cite{ProveIt2026}.

The relevant preprocessor is in
\code{Topology/UnknotRecognition/fast/fastunknot/simplify.py}, reviewed blob
\hashcode{2fcaabb51eaa45cf84bdee32cd1b4e18acf2c478}.  The file represents a diagram
by a fixed list of crossing darts and a mutable edge involution.  It performs
crossing-decreasing RI/RII moves incrementally.  When no such move remains, it
can trial an RIII inversion of a triangular face, keep a short sequence if the
sequence exposes RI/RII, and otherwise undo it exactly.

The existing recursive helper has three important properties.
\begin{enumerate}
\item It is \emph{sound}: successful RIII trials are retained only when a legal
      RI/RII move appears, and the resulting trace is independently replayable.
\item It is \emph{budgeted}: the total number of trial moves is capped per input
      crossing, because failed deep searches can dominate runtime.
\item It is \emph{path-local}: after the first RIII move, later candidates are
      drawn from a neighborhood of the immediately preceding move, rather than
      from the accumulated support of the entire trace.
\end{enumerate}

The synthesis reports document why this matters.  On a corpus of difficult
unknot survivors, depth two, three, and four searches simplified many more
instances than one-step lookahead, but a failed depth-four search on a
triangle-rich braid was much more expensive.  The current default therefore
uses shallow depth and a strict budget.  These data motivate a theorem that
separates global branching from genuinely local continuation.

\subsection{Relation to known Reidemeister bounds}

Hass and Lagarias proved an exponential upper bound on the number of
Reidemeister moves needed to trivialize an $n$-crossing unknot diagram
\cite{HassLagarias2001}.  Dynnikov proved monotonic simplification for
arc-presentations of the unknot \cite{Dynnikov2006}; Henrich and Kauffman used
that framework to bound diagram complexity and move counts
\cite{HenrichKauffman2011}.  Lackenby later proved that an $n$-crossing unknot
diagram can be reduced using at most $(236n)^{11}$ Reidemeister moves while the
crossing number remains at most $(7n)^2$ along the sequence
\cite{Lackenby2015}.

Those theorems establish polynomial-size certificates and powerful global
control.  They do not assert that a sequence can be partitioned into bounded
runs of RIII moves separated only by crossing-\emph{decreasing} RI/RII moves.
The present search deliberately forbids crossing increase.  Consequently, the
new parameterized theorem is neither a replacement for Lackenby's theorem nor a
consequence of it.

Lackenby's 2026 hierarchy algorithm gives a different route to unknot detection,
through incompressibility and essential hierarchies \cite{Lackenby2026}.  Its
abstract emphasizes avoidance of certain large searches, but the preprint does
not claim the quasi-polynomial running-time theorem sought by this project.  The
local RIII analysis here is orthogonal and can coexist with hierarchy-based
certificates.

\subsection{Relation to the exact fallback}

Khovanov homology detects the unknot \cite{KronheimerMrowka2011}, and Bar--Natan's
tangle/cobordism formalism supports local scanning computations
\cite{BarNatan2005}.  ProveIt exploits this exact route after cheaper stages
fail.  Report 25 changes none of that algebra.  Its role is to make one
preprocessing search asymptotically accountable and more complete under explicit
parameters.

\subsection{Methodological scope}

The OpenAI mathematics repository contains manuscripts and proof artifacts at
multiple verification levels, explicitly warning that unformalized outputs may
contain errors \cite{OpenAIMath2026}.  This report follows the useful part of
that model: theorem statements are separated from experiments, failed proof
ideas are recorded, and the integration patch remains opt-in.  It does not rely
on repository size or model provenance as mathematical evidence.

\begin{warning}[Novelty and verification]
The theorems below are new relative to the reviewed ProveIt source and its
synthesis, but no claim of priority in the wider mathematical literature is made
without expert review.  The proofs are human-readable and computationally
stress-tested, not formally verified.  The ProveIt adapter was not executed
inside a full checkout in this environment.
\end{warning}

\section{A bounded-local reversible rewrite model}
\label{sec:model}

This section isolates exactly the hypotheses used later.  The model is broad
enough for the fixed-dart Reidemeister implementation and narrow enough that
commutation arguments are valid without hidden geometric assumptions.

\subsection{States, actions, and stable sites}

Fix a finite set $V$ of \emph{sites}, with $|V|\leq c_0N$ for an input-size
parameter $N$.  A state belongs to a set $\Omega_N$.  There are two kinds of
local operation.

\begin{definition}[Neutral action]
A neutral action instance $a$ legal at $x\in\Omega_N$ is a reversible transition
\[
  x \xrightarrow{a} a x
\]
with a footprint $\Phi_x(a)\subseteq V$.  We assume
$|\Phi_x(a)|\leq f$, where $f$ is independent of $N$.
\end{definition}

\begin{definition}[Terminal reduction]
A terminal reduction instance $r$ legal at $x$ is a local decreasing operation
with footprint $\Psi_x(r)\subseteq V$, with
$|\Psi_x(r)|\leq g$.  During an unlocking search we only ask whether $r$ is
legal; the reduction is applied after a trace succeeds.
\end{definition}

For RIII, the neutral action preserves the number of crossings.  RI/RII are the
terminal reductions.  Stable sites are fixed input darts, so a crossing can be
moved combinatorially without renaming its darts.

\subsection{Strong locality}

The main axiom is a two-sided commutation diamond.  It is stated for executed
action instances, because a later move need not be syntactically visible before
an earlier local update.

\begin{definition}[Strong disjoint commutation]
A rewrite system has strong disjoint commutation when the following hold.
\begin{enumerate}[label=(L\arabic*)]
\item \textbf{Backward and forward diamond.}  Suppose
      $x\xrightarrow{a}y\xrightarrow{b}z$, and the residual footprints of $a$
      and $b$ are disjoint.  Then $b$ is legal at $x$, $a$ is legal after $b$,
      both orders reach the same state, and the two footprints are unchanged by
      swapping the order.
\item \textbf{Terminal locality.}  If
      $x\xrightarrow{a}y$ and a terminal reduction $r$ is legal at one endpoint
      with footprint disjoint from $\Phi_x(a)$, then its corresponding residual
      is legal at the other endpoint with the same footprint.
\item \textbf{Reversibility.}  Every neutral action has a local inverse whose
      footprint is the same bounded set of sites.
\end{enumerate}
\end{definition}

The backward half of (L1) is essential.  It says that a move observed after an
unrelated action was already legal before that action.  This is what permits a
birth action to commute left.

\begin{remark}
For a concrete implementation, one can use a conservative footprint larger than
the minimal topological support.  This can only classify more actions as
interacting.  The search becomes less sharply pruned, but all commutation claims
remain safe.
\end{remark}

\subsection{Incidence bounds}

We separate the number of moves available globally from the number meeting a
small active region.

\begin{definition}[Bounded incidence]
A rule family has constants $A,B>0$ when, for every state $x$,
\begin{align*}
  |\{a:a\text{ neutral and legal at }x\}| &\leq AN, \\
  |\{a:a\text{ neutral and legal at }x,
      \ \Phi_x(a)\cap S\neq\varnothing\}| &\leq B|S|
\end{align*}
for every $S\subseteq V$.
\end{definition}

The first bound permits $O(N)$ possible locations for a new independent branch.
The second says that continuation from a region of $O(k)$ sites has only $O(k)$
choices.  In a finite local rule system on a bounded-degree combinatorial
structure, such constants are natural.

\subsection{Unlocking traces}

Let
\[
 x_0\xrightarrow{a_1}x_1\xrightarrow{a_2}\cdots
 \xrightarrow{a_k}x_k
\]
be a neutral trace.

\begin{definition}[First-unlocking trace]
The trace is first-unlocking when no terminal reduction is legal at
$x_0,\ldots,x_{k-1}$ and at least one terminal reduction $r$ is legal at $x_k$.
Its length is $k$.  It is \emph{shortest} when no first-unlocking trace from
$x_0$ has smaller length.
\end{definition}

A search can stop at the first terminal reduction.  Therefore any trace that
exposes a reduction earlier than expected is automatically replaced by its
shorter prefix.

\begin{definition}[Birth and birth number]
Write $\Phi_i=\Phi_{x_{i-1}}(a_i)$.  Action $a_i$ is a \emph{birth} when
\[
  \Phi_i\cap\bigcup_{j<i}\Phi_j=\varnothing.
\]
The birth number $\beta(\tau)$ is the number of birth actions in $\tau$.
\end{definition}

The first action is always a birth.  A birth is not necessarily topologically
remote forever: later actions may connect its footprint to other branches.

\subsection{The action-support graph}

Fix a terminal reduction $r$ legal at $x_k$.  The action-support graph
$H(\tau,r)$ has vertices $1,\ldots,k,\star$.  Vertex $i$ represents $a_i$ and
$\star$ represents $r$.  Put an edge between two action vertices when their
executed footprints intersect, and an edge $i\star$ when $\Phi_i$ intersects
the terminal footprint $\Psi_{x_k}(r)$ after the canonical residual
identification supplied by strong locality.

In the fixed-site situations used here, the residual identification is literal:
the same dart labels are compared throughout.  The graph records causal
noncommutation, not chronological adjacency.

\begin{example}[Two branches before a join]
Let actions $A$ and $B$ have disjoint footprints.  Let $C$ overlap both, and let
the terminal reduction overlap $C$.  The chronological word $ABC$ has two
support components after its second letter, but the full action-support graph is
connected.  This is the smallest obstruction to the naive claim that every
shortest trace can be given connected prefixes.
\end{example}

\section{Causal structure and the birth-front normal form}
\label{sec:causal}

\subsection{Shortest traces have connected dependency graphs}

\begin{theorem}[Causal connectedness]
\label{thm:causal-connected}
Assume strong disjoint commutation.  If $\tau$ is a shortest first-unlocking
trace and $r$ is a terminal reduction exposed at its end, then
$H(\tau,r)$ is connected.
\end{theorem}

\begin{proof}
Suppose $H(\tau,r)$ is disconnected.  Let $C$ be a connected component not
containing the terminal vertex $\star$.  Every action in $C$ has footprint
disjoint from every action outside $C$ and from the terminal reduction.

Use the commutation diamond repeatedly to move every action of $C$ to the right,
preserving the relative order inside $C$ and the relative order of all actions
outside $C$.  Each adjacent swap is legal, preserves the reached state, and
preserves the relevant footprints.  We obtain a trace-equivalent word
\[
   \tau'\,\sigma_C,
\]
where $\sigma_C$ consists exactly of the actions in $C$.

The terminal reduction is legal after $\tau'\sigma_C$.  Since its footprint is
disjoint from every footprint in $\sigma_C$, terminal locality moves it left
across the entire block.  Hence it is already legal after $\tau'$.  Deleting
$\sigma_C$ gives a strictly shorter trace that exposes a terminal reduction;
if one was exposed at an earlier prefix of $\tau'$, truncate at its first such
prefix.  In either case we obtain a strictly shorter first-unlocking trace from
$x_0$, contradicting the minimality of $\tau$.
\end{proof}

\begin{corollary}
Every birth in a shortest trace has a path, through overlapping later action
footprints, to the terminal reduction.
\end{corollary}

This corollary is qualitative.  It rules out irrelevant disconnected work, but
it does not bound the number of births.  A connected dependency tree can have
many leaves.

\subsection{Why connected prefixes are false}

The following proposition records the proof-audit failure that led to the birth
parameter.

\begin{proposition}[Branching obstruction]
\label{prop:branching-obstruction}
There exists a strongly local reversible rewrite system with a shortest
first-unlocking trace of length three such that every first-unlocking trace of
length three has birth number two.  In particular, no commuting-equivalent
first-unlocking trace has connected chronological prefixes.
\end{proposition}

\begin{proof}
Use three Boolean sites $0,1,2$, initially all zero.  Action $A$ toggles site
$0$ and has footprint $\{0\}$.  Action $B$ toggles site $2$ and has footprint
$\{2\}$.  Action $C$ is legal exactly when sites $0$ and $2$ are one; it toggles
site $1$ and has footprint $\{0,1,2\}$.  The terminal reduction is legal exactly
when site $1$ is one.

The moves are involutions, and disjoint supports commute.  Neither $A$ nor $B$
alone enables $C$, so at least two moves are required before $C$.  The words
$ABC$ and $BAC$ unlock the reduction in three moves.  Their first two footprints
are disjoint, so each has two births.  No trace with one birth can set both sites
$0$ and $2$ before applying $C$.
\end{proof}

The fixture is implemented as \code{branching\_example()} in the prototype.  It
is not asserted to be a knot diagram; its purpose is to refute an invalid
general lemma before that lemma is imported into knot topology.

\subsection{Births commute to the front}

\begin{lemma}[Births are pairwise disjoint]
\label{lem:birth-disjoint}
The footprints of distinct births in one trace are disjoint.
\end{lemma}

\begin{proof}
A later birth is disjoint from the union of every earlier footprint, including
the footprint of any earlier birth.
\end{proof}

\begin{lemma}[Left commutation of a birth]
\label{lem:birth-left}
A birth action can be commuted left across every action preceding it.
\end{lemma}

\begin{proof}
By definition, its footprint is disjoint from each preceding footprint.  Apply
the backward-and-forward commutation diamond one adjacent predecessor at a
time.  Footprints remain unchanged under each swap, so the next swap remains
valid.
\end{proof}

\begin{theorem}[Birth-front normal form]
\label{thm:birth-front}
Let $\tau$ be a first-unlocking trace of length $k$ and birth number $b$ in a
strongly local reversible system.  There is a first-unlocking trace
$\widehat\tau$ of length at most $k$ and birth number at most $b$ such that:
\begin{enumerate}
\item all births of $\widehat\tau$ occur before every nonbirth action;
\item after the birth block, every action footprint intersects the union of all
      earlier footprints.
\end{enumerate}
\end{theorem}

\begin{proof}
Identify the births in the original trace.  By Lemma~\ref{lem:birth-left}, move
each birth left across all predecessors.  By
Lemma~\ref{lem:birth-disjoint}, births commute with one another, so choose any
fixed canonical order for them.  Preserve the relative order of all nonbirth
actions.

Consider an original nonbirth action $a_i$.  It overlapped at least one earlier
action.  If that earlier action was a birth, it has moved farther left.  If it
was a nonbirth, their relative order was preserved.  Thus $a_i$ still meets the
accumulated earlier footprint in the reordered word.  Strong commutation shows
that the reordered word is legal and reaches the same final state.

It is possible that the terminal reduction becomes legal at an earlier prefix
of the reordered word.  In that case truncate at the first such prefix.  The
length and number of births can only decrease.  The truncated word is
first-unlocking and has the two stated properties.
\end{proof}

\begin{remark}[Normal form versus a chronological invariant]
The birth number is defined on a trace, not just its final state.  The theorem
shows that it has algorithmic meaning: it is an upper bound on the number of
global choices needed before all remaining choices are local.  One may minimize
$b$ over shortest traces to obtain an input parameter, but the search does not
need to know that minimum in advance; it can deepen in $b$.
\end{remark}

\subsection{Accumulated support strictly improves last-touch search}

The current ProveIt recursion follows the neighborhood of the immediately
previous RIII move.  The next proposition shows that even the one-birth normal
form is strictly more general in local rewrite systems.

\begin{proposition}[Strict separation]
\label{prop:last-vs-accumulated}
There is a strongly local reversible system with a one-birth first-unlocking
trace of length three, but no first-unlocking trace of length at most three in
which consecutive action footprints always intersect.
\end{proposition}

\begin{proof}
Use sites $0,1,2,3$, initially zero.  Action $A$ toggles sites $0$ and $1$ and
has footprint $\{0,1\}$.  Action $B$ requires site $0=1$, toggles site $2$, and
has footprint $\{0,2\}$.  Action $C$ requires site $1=1$, toggles site $3$, and
has footprint $\{1,3\}$.  The terminal reduction requires sites $2=3=1$.

After $A$, both $B$ and $C$ are legal.  They commute because their footprints
are disjoint.  Either $ABC$ or $ACB$ unlocks the terminal reduction.  There is
one birth, namely $A$, and each leaf meets the accumulated footprint of $A$.
However, $B$ and $C$ do not meet one another, so neither length-three ordering
has pairwise-intersecting consecutive footprints.  No shorter word can set both
sites $2$ and $3$.
\end{proof}

The fixture is \code{accumulated\_strict\_example()} in the prototype.  For knot
diagrams, this proposition motivates the accumulated-footprint search; it does
not by itself exhibit a concrete RIII diagram with the same separation.

\section{Parameterized enumeration and asymptotic regimes}
\label{sec:enumeration}

\subsection{The birth-front search}

The normal form leads to a direct iterative-deepening algorithm.

\begin{algorithmic}[Birth-front unlocking search]
\label{alg:birth-front}
Given a state $x$, a depth bound $k$, and a birth bound $b$:
\begin{enumerate}
\item For $\ell=1,\ldots,k$, search traces of length at most $\ell$.
\item Enumerate a canonical list of at most $b$ pairwise footprint-disjoint
      births.  Each birth is chosen among all currently legal neutral actions.
\item After the birth block, permit only a legal action whose footprint meets
      the union of all previously chosen footprints.
\item After every neutral action, test for a terminal reduction.  Return the
      shortest successful trace.
\item Undo failed branches exactly.  Exhaustion returns \emph{unknown}, never a
      negative verdict.
\end{enumerate}
\end{algorithmic}

A canonical order on births removes permutations of pairwise commuting actions.
The complexity theorem below does not need that saving, so its bound is robust
to a simpler implementation.

\subsection{Counting search nodes}

\begin{theorem}[Parameterized enumeration]
\label{thm:enumeration}
Assume footprint size at most $f$, global incidence at most $AN$, local
incidence at most $B|S|$, and polynomial-time state operations.  The birth-front
algorithm finds a first-unlocking trace whenever one of length at most $k$ and
birth number at most $b$ exists, and exhausts all birth-front normal-form
candidates in time
\[
   N^{b+O(1)}(Ck)^k,
   \qquad C=Bf\max(1,A),
\]
and polynomial working memory.  The constant implicit in $O(1)$ is independent
of $b$ and $k$.
\end{theorem}

\begin{proof}
By Theorem~\ref{thm:birth-front}, it is enough to enumerate traces with all
births first.  Suppose a normal-form trace has $r\leq b$ births and total length
$\ell\leq k$.

At each birth position there are at most $AN$ legal neutral actions, so the birth
block has at most $(AN)^r$ possibilities.  After $j$ actions have been chosen,
the accumulated footprint contains at most $fj$ sites.  Bounded local incidence
therefore gives at most $Bfj$ possible connected extensions.  Consequently,
the number of length-$\ell$ search words with $r$ births is at most
\[
  (AN)^r\prod_{j=r}^{\ell-1} Bfj
  \leq (AN)^b(Bfk)^k.
\]
Summing over $r\leq b$ and $\ell\leq k$ contributes at most a polynomial or
$2^{O(k)}$ factor, which is absorbed by replacing $C$ with a larger absolute
rule-set constant.

Each node performs polynomial-time legality, application, terminal, and undo
operations.  A depth-first traversal stores one state representation, at most
$k$ undo records, the accumulated footprint of size at most $fk$, and polynomial
bookkeeping.  Hence the working memory is polynomial in $N$ and $k$.
\end{proof}

\begin{remark}[Why the exponent $b$ appears]
Every birth can be located anywhere in the diagram, so a direct deterministic
search pays one factor of $N$ per independent branch.  Removing this exponent
would require extra structure, such as color-coding, representative families,
a reverse terminal-rooted search, or a geometric theorem constraining where
births can occur.
\end{remark}

\subsection{Quasi-polynomial and polynomial corollaries}

Write
\[
 T(N;b,k)=N^{b+O(1)}(Ck)^k.
\]
Taking logarithms gives
\[
 \log T=(b+O(1))\log N + k\log(Ck).
\]

\begin{corollary}[Logarithmic parameters]
\label{cor:log-parameters}
If $b=O(\log N)$ and $k=O(\log N)$, then
\[
 T(N;b,k)=\exp(O((\log N)^2))=N^{O(\log N)}.
\]
Thus the search is quasi-polynomial.
\end{corollary}

\begin{corollary}[Constant branching]
\label{cor:constant-branching}
If $b=O(1)$ and $k=O(\log N)$, then
\[
 T(N;b,k)=\exp(O(\log N\log\log N))
          =N^{O(\log\log N)}.
\]
\end{corollary}

\begin{corollary}[Sublogarithmic depth]
\label{cor:polynomial-island}
If $b=O(1)$ and
$k=O(\log N/\log\log N)$, then $T(N;b,k)=N^{O(1)}$.
\end{corollary}

\begin{proof}
For the last statement,
$k\log(Ck)=O(\log N)$, while the birth factor is polynomial because $b$ is
constant.
\end{proof}

\subsection{Search budgets}

A practical implementation will cap the number of trial moves rather than run
the worst-case tree.  The theorem and the budget serve different purposes.

\begin{itemize}
\item The theorem says what is complete when the specified $(b,k)$ tree is
      exhausted.
\item A trial budget may stop earlier to protect the exact fallback.
\item Budget exhaustion has no mathematical meaning about knottedness or about
      the nonexistence of a trace.
\item Successful traces remain exact because they are independently replayed.
\end{itemize}

This is compatible with the current ProveIt design, which already treats RIII
search as an optional pre-scan accelerator.

\section{Specialization to Reidemeister-III unlocking}
\label{sec:r3}

We now verify the abstract hypotheses for the mutable dart representation used
by the reviewed ProveIt code.  The purpose is not to reprove Reidemeister's
theorem; it is to account for the local search tree.

\subsection{Fixed darts and the face permutation}

An $N$-crossing diagram has darts
\[
  V=\{0,1,\ldots,4N-1\}.
\]
Darts $4i,\ldots,4i+3$ belong to crossing $i$.  Let $\alpha$ be the fixed-point
free involution pairing the two darts of each diagram edge.  Let $\rho$ rotate a
dart to the next slot at its crossing.  A face walk is generated by
\[
  \varphi=\rho\circ\alpha.
\]
The implementation writes this map as \code{nxt}.  A triangular face is a
three-cycle of $\varphi$ through three distinct crossings.  The over/under slot
parities determine whether that triangle is an admissible RIII configuration.

The RIII update keeps crossing identities and slots fixed and modifies only the
edge involution.  This is particularly convenient for the causal analysis:
site names do not move even though adjacency does.

\subsection{A conservative RIII footprint}

Let $T$ be a legal RIII triangle with crossing set $C(T)$ of size three.  Define
\[
  D(T)=\{4c+j:c\in C(T),\ 0\leq j<4\}
\]
and the conservative footprint
\[
  \Phi(T)=D(T)\cup\alpha(D(T)).
\]
Thus $|\Phi(T)|\leq24$.

\begin{lemma}[Footprint containment]
\label{lem:r3-footprint}
The reviewed \code{apply\_r3} routine changes $\alpha$ only on darts in
$\Phi(T)$.  Its inverse has the same footprint set.
\end{lemma}

\begin{proof}
The routine names the six outer slots of the three participating strands and
the corresponding inner slots.  Every new link has one endpoint at a dart of
$D(T)$; its other endpoint is either another dart of $D(T)$ or the previous
$\alpha$-partner of a dart in $D(T)$.  Hence every written array entry lies in
$D(T)\cup\alpha(D(T))$.  Moreover, each outside dart in
$\alpha(D(T))\setminus D(T)$ remains paired to some dart of $D(T)$ after the
move, while no new outside dart is introduced.  Therefore the post-move set
$D(T)\cup\alpha'(D(T))$ equals the pre-move footprint.  RIII is an involution,
so the reverse action has exactly this same footprint.
\end{proof}

For RI/RII terminal tests, use the analogous conservative footprint consisting
of all darts at the one or two participating crossings together with their edge
partners.  Its size is bounded by 16.

\subsection{Disjoint footprints commute}

\begin{proposition}[Strong locality for the dart rules]
\label{prop:r3-locality}
RIII actions and RI/RII terminal tests in the fixed-dart implementation satisfy
strong disjoint commutation when the conservative footprints above are used.
\end{proposition}

\begin{proof}
An RIII action writes only entries of $\alpha$ in its footprint by
Lemma~\ref{lem:r3-footprint}.  Suppose a second RIII action is observed after
the first and their residual footprints are disjoint.  The twelve darts owned
by the second crossing triple lie outside the first footprint.  Since the first
move changes $\alpha$ only inside its own footprint, the partners of those
twelve darts are also unchanged; hence the second conservative footprint, its
three-cycle face test, and its slot-parity test were already identical before
the first move.  The same argument in the other direction proves the two-sided
legality diamond.  The two updates write disjoint array entries, so they commute
entry by entry and their footprints remain fixed under the swap.

The same argument applies to RI/RII legality: monogon and bigon face walks read
only the darts at their participating crossings and their partners.  A
disjoint RIII update neither creates nor destroys such a local face.  Equivalently,
any RI/RII move newly exposed by an RIII update has a conservative footprint
intersecting that RIII footprint.  Finally, RIII inversion is its own local
inverse with the same footprint.
\end{proof}

\begin{remark}
Two RIII moves with disjoint \emph{crossing triples} need not have disjoint
conservative footprints, because their local disks can touch the same outside
edge.  This is why the theorem is formulated with dart footprints rather than
only crossing labels.
\end{remark}

\subsection{Global and local incidence}

For a connected spherical $4$-regular projection graph with $N$ vertices,
$E=2N$ and Euler's formula gives $F=N+2$ faces.  A legal RIII action is supported
by a triangular face, so the total number of legal RIII locations is at most
$N+2$.  This gives the global incidence bound with an absolute $A$.

For the local bound, fix one dart $d$.  If $d\in\Phi(T)$, then either a crossing
of $T$ owns $d$, or a crossing of $T$ owns $\alpha(d)$.  Each crossing has four
face corners, so it is incident to at most four triangular faces.  Consequently
at most eight triangular faces have footprints containing a given dart.  By a
union bound,
\[
  |\{T:\Phi(T)\cap S\neq\varnothing\}|\leq 8|S|.
\]
Thus one may take $f=24$ and $B=8$, giving $Bf=192$.  Whenever an RIII
location exists, $N\geq3$ and $N+2\leq2N$, so a canonical enumeration with one
representative per triangular face may take $A=2$ and hence $C=384$ in
Theorem~\ref{thm:enumeration}.  An implementation that emits orientations or
duplicate representatives is covered by some larger absolute $C$; no numerical
optimality is claimed, and the asymptotic regimes are unchanged.

\begin{corollary}[RIII search bound]
\label{cor:r3-bound}
On an $N$-crossing fixed-dart state, a first-unlocking RIII trace of length at most $k$ and birth number at most
$b$, when one exists, can be found by exhaustive birth-front search in
\[
  N^{b+O(1)}(Ck)^k
\]
time and polynomial memory for an absolute constant $C$.
\end{corollary}

\subsection{Why accumulated footprints are the right implementation object}

The existing recursive search computes a neighborhood from the triangle just
applied and recursively scans that neighborhood.  This is a last-touch chain.
For birth number one, the complete normal-form search must instead retain the
union of all previous RIII footprints.  A candidate is local when its current
footprint intersects that union.  The active set contains at most $24k$ darts,
so scanning crossings owning active darts and their current partners generates
only $O(k)$ local triangle candidates, up to duplicates.

For $b>1$, all births are chosen first.  They are required to have pairwise
disjoint footprints and are ordered canonically.  After the first connected
move, no further birth is permitted.  This exactly implements
Theorem~\ref{thm:birth-front}.

\subsection{Interaction with the face shortcut}

The current implementation's \code{r3\_can\_help} test is used only for the
last allowed RIII move.  It checks a necessary local condition for that move to
create an RI/RII face.  Retaining that placement is important.  Applying the
shortcut to every intermediate move would be unsound, because an intermediate
RIII need not immediately expose a reduction; it may only prepare a later
interaction.

\subsection{State restoration}

Every unsuccessful branch must restore the exact involution.  The reviewed code
undoes an RIII move by locating the inverse triangle and applying the same
involution routine.  The adapter preserves this contract.  A production
integration should additionally test every budget-exhaustion boundary, because
an early return from a recursive branch is a common place for mutable-state
leaks.

\section{The restricted recognizer and integration design}
\label{sec:algorithm}

\subsection{A monotone simplification class}

Fix functions $b(N)$ and $k(N)$, a deterministic RI/RII exhaustion order,
and a deterministic ordering of the birth-front search tree (for example,
iterative depth followed by lexicographic triangle keys).  Let
$\mathsf{BF}_{b,k}$ denote the resulting deterministic preprocessor.  Define
$\mathcal{U}^{\mathsf{BF}}_{b,k}$ to be the set of unknot diagrams on which this
specific process reaches the crossing-free circle:
\begin{enumerate}
\item exhaust RI/RII in the fixed order;
\item if crossings remain, run the complete birth-front search with bounds
      $b(N),k(N)$ and select its first trace in the fixed ordering;
\item apply that trace, return to step 1, and return $\Unknown$ if no trace is
      found.
\end{enumerate}
The policy qualification is necessary.  Two different successful unlock traces
can lead to different reduced diagrams, and the existence of one globally good
sequence does not by itself justify a greedy choice without backtracking.  The
definition permits the diagram size to fall during the process; using the
original $N$ as the upper-bound parameter only weakens the estimate.

\begin{theorem}[Restricted-class unknot recognition]
\label{thm:restricted-recognition}
There is a deterministic, sound algorithm that returns $\Unk$ on every diagram
in $\mathcal{U}^{\mathsf{BF}}_{b,k}$ and otherwise returns either $\Unk$ or $\Unknown$.  Its
running time is
\[
  N^{b(N)+O(1)}(Ck(N))^{k(N)}
\]
and its working memory is polynomial, apart from any subsequent fallback.
\end{theorem}

\begin{proof}
At each stalled state, run the complete birth-front search with the specified
bounds and fixed ordering.  If it succeeds, select the first trace in that
ordering, replay it, expose RI/RII, and exhaust reductions in the fixed order.  Reidemeister moves preserve the represented
knot type.  Therefore reaching the crossing-free circle is a sound unknot
certificate.

Each successful episode is followed by an RI or RII move and hence removes at
least one crossing.  There are at most $N$ episodes.  Multiplying the
single-episode bound from Corollary~\ref{cor:r3-bound} by $N$ is absorbed in the
$N^{O(1)}$ factor.  If a search fails, return $\Unknown$ or invoke an exact
fallback; no negative conclusion is drawn.
\end{proof}

\begin{corollary}
The policy class $\mathcal{U}^{\mathsf{BF}}_{b,k}$ is recognized in quasi-polynomial time when
$b(N),k(N)=O(\log N)$, in $N^{O(\log\log N)}$ time when
$b(N)=O(1)$ and $k(N)=O(\log N)$, and in polynomial time when
$b(N)=O(1)$ and $k(N)=O(\log N/\log\log N)$.
\end{corollary}

\subsection{Certificate size and replay}

A successful preprocessor emits the same kind of trace already used by ProveIt:
for each RIII move, the three stable input crossing indices, followed by RI/RII
records.  Across at most $N$ reduction episodes, a uniform depth bound $k$ gives
$O(Nk)$ RIII records.  Each crossing index uses $O(\log N)$ bits, so the
certificate size is $O(Nk\log N)$ bits.  The existing independent replay routine
checks legality move by move and rebuilds a validated spherical one-component
diagram.

This proof-carrying behavior is more important than the heuristic that found the
trace.  Search policies, budgets, and branch ordering may change without
changing the trusted replay kernel.

\subsection{Opt-in API}

The supplied review patch adds two optional arguments to
\code{simplify}:
\begin{center}
\begin{tabular}{@{}lll@{}}
\toprule
Argument & Suggested values & Meaning \\
\midrule
\code{r3\_search} & \code{"last"}, \code{"clustered"} & legacy or birth-front search \\
\code{r3\_births} & positive integer & maximum number of births \\
\bottomrule
\end{tabular}
\end{center}
The default remains \code{"last"}; therefore applying the patch alone does not
change existing behavior.  In clustered mode, \code{r3\_births=1} gives the
accumulated-support search.  Larger values admit independent birth branches.
The existing \code{r3\_depth}, \code{r3\_budget}, and face-shortcut controls are
retained.

\subsection{Pseudocode}

A simplified version of one depth-limited search is shown below.

\begin{lstlisting}[language=Python,caption={Birth-front RIII search, schematic}]
def search(state, depth_left, births_left, birth_phase,
           active_footprint, last_birth_key, path):
    if budget_exhausted() or depth_left == 0:
        return failure

    if birth_phase and births_left > 0:
        for T in all_legal_triangles(state):
            F = footprint(state, T)
            if F intersects active_footprint:
                continue
            if key(T) <= last_birth_key:
                continue
            apply(T)
            if RI_or_RII_now_legal(F):
                return success
            if search(state, depth_left - 1, births_left - 1,
                      True, active_footprint union F, key(T), path + [T]):
                return success
            undo(T)

    for T in legal_triangles_meeting(state, active_footprint):
        F = footprint(state, T)
        apply(T)
        if RI_or_RII_now_legal(F):
            return success
        if search(state, depth_left - 1, births_left,
                  False, active_footprint union F,
                  last_birth_key, path + [T]):
            return success
        undo(T)

    return failure
\end{lstlisting}

The production helper contains duplicate suppression, inverse-move avoidance,
trial accounting, and the last-move face shortcut.

\subsection{Required production tests}

Before clustered mode can be enabled by default, the following checks are
mandatory.

\begin{enumerate}
\item \textbf{Trace replay.}  Replay every successful trace with the existing
      independent checker.
\item \textbf{Exact restoration.}  Hash the entire $\alpha$ involution before
      and after every failed search and every budget cutoff.
\item \textbf{Invariant preservation.}  Compare Alexander data, writhe, and
      Khovanov ranks on randomized successful RIII traces, as the existing
      synthesis did for the legacy helper.
\item \textbf{Verdict equality.}  Compare the full recognizer with clustered
      mode disabled and enabled on maintained and random corpora.
\item \textbf{Adversarial branching.}  Add fixtures requiring two disjoint
      birth footprints before a connecting RIII move, if a genuine knot-diagram
      example can be found.
\item \textbf{Performance separation.}  Measure successes and failures
      separately.  Deep failed searches should remain budgeted and should run
      only immediately before the expensive exact scan.
\end{enumerate}

\subsection{Why the exact fallback remains essential}

Even a complete $(b,k)$ search is incomplete for unrestricted unknot diagrams.
If it fails, the diagram may still be an unknot that needs larger parameters,
crossing-increasing moves, a different diagram language, or a non-Reidemeister
certificate.  The ProveIt pipeline should therefore keep its exact Khovanov
recognizer and every sound obstruction.  Clustered RIII search is a front-end
optimization with a theorem, not a replacement decision procedure.

\section{Executable validation and negative controls}
\label{sec:validation}

The mathematical proof does not depend on experiments.  The experiments serve
three narrower purposes: catch implementation mistakes, demonstrate the two
strict separations, and make the scope of evidence explicit.

\subsection{Reference implementation}

The file \code{prototype/causal\_search.py} implements a finite class of Boolean
local rewrite systems.  A neutral action has:
\begin{itemize}
\item a stable support;
\item guards that read only support sites;
\item toggles on support sites not read by the guards.
\end{itemize}
Applying a legal action twice restores the state.  Actions with disjoint
supports commute and preserve one another's legality.  A terminal reduction is
a local bit-pattern predicate.

Three searches are included.
\begin{enumerate}
\item unrestricted breadth-first search, used as a finite oracle;
\item the birth-front search proved complete under the parameter bound;
\item a last-touch search modeling the current immediate-predecessor locality
      restriction.
\end{enumerate}

\subsection{Hand-checked strict examples}

The two examples from Propositions~\ref{prop:branching-obstruction} and
\ref{prop:last-vs-accumulated} are executable.  A third fixture checks a
subtler implementation point: accumulated support is search state even when the
physical rewrite state repeats.

\begin{center}
\begin{tabular}{@{}lccc@{}}
\toprule
Fixture & Witness depth & Minimum births & Distinguished failure \\
\midrule
Two independent births, then join & 3 & 2 & one-birth search fails \\
One birth, two commuting leaves & 3 & 1 & last-touch search fails \\
State repeat extends support & 4 & 1 & state-only pruning fails \\
\bottomrule
\end{tabular}
\end{center}

These are negative controls.  They ensure the implementation is not silently
reinstating either rejected claim.

\subsection{Generated finite systems}

The deterministic validation script generated 5,000 systems with six Boolean
sites, nine neutral actions, three terminal reductions, and search depth at most
six.  The random seed is the system index.  The outcome was:

\begin{center}
\begin{tabular}{@{}lr@{}}
\toprule
Category & Count \\
\midrule
Initially reducible & 3,209 \\
No witness within depth six & 131 \\
Witness found by unrestricted BFS & 1,660 \\
Front-loading failures & 0 \\
Birth-front completeness discrepancies & 0 \\
Disconnected shortest dependency graphs & 0 \\
\bottomrule
\end{tabular}
\end{center}

Among the 1,660 witnessed systems, the unrestricted oracle returned depths
$1,2,3,4,5$ with counts $1381,223,47,8,1$.  The birth-front search needed birth
bound one in 1,603 cases and birth bound two in 57 cases.  The oracle's first
returned shortest trace had birth counts one, two, and three in 1,548, 110, and
two cases respectively; a different shortest trace sometimes used fewer births.

The unrestricted BFS explored 3,863 dequeued states across these runs; the
iterative birth-front comparisons explored 12,457 recursive nodes.  This is not
a speed benchmark: the state space has only 64 states, BFS deduplicates globally,
and the parameterized search deliberately mirrors a local depth-first
implementation.  The relevant result is equality of witnessed behavior, not a
runtime win on toy Boolean systems.

\subsection{Exhaustive three-site catalog}

A second script exhausts a fully specified finite catalog rather than sampling.
There are 27 one-bit involution templates on three sites: choose the toggled
site, and for each of the other two sites choose guard value zero, guard value
one, or absence.  The script checks every three-move subset, every one of the 26
nonempty local terminal predicates, and all eight initial states, for a total of
608,400 instances through depth four.

\begin{center}
\begin{tabular}{@{}lr@{}}
\toprule
Catalog outcome & Count \\
\midrule
Initially reducible & 163,800 \\
No witness within depth four & 272,632 \\
Witnessed & 171,968 \\
Front-loading failures & 0 \\
Birth-front completeness discrepancies & 0 \\
Disconnected shortest dependency graphs & 0 \\
\bottomrule
\end{tabular}
\end{center}

This catalog is still only a finite software check, but it removes sampling
uncertainty inside the stated model class.  The exact parameters and counts are
stored in \code{results/exhaustive\_catalog.json}.

\subsection{Unit tests}

Five generic-search and two adapter-contract checks cover:
\begin{enumerate}
\item the two-birth branching fixture;
\item strict improvement of accumulated support over last-touch locality;
\item preservation or shortening under front-loading;
\item the support-history state-repeat regression;
\item 200 additional generated-system oracle comparisons;
\item exact adapter restoration after failure; and
\item footprint-stable inverse replay after adapter success.
\end{enumerate}
All seven checks pass under the Python interpreter in the artifact environment.

\subsection{What was not tested}

No full ProveIt checkout was available to the container runtime.  Consequently:
\begin{itemize}
\item the 431 maintained repository tests were not rerun;
\item the adapter was exercised only against a twelve-dart contract gadget,
      not actual validated PD diagrams;
\item no timing claim is made for clustered RIII search;
\item no knot fixture is claimed to separate last-touch and accumulated support;
\item no corpus claim is extrapolated from the Boolean model.
\end{itemize}

The integration snippet was syntax-checked and its success/failure mutation
contract was tested against a twelve-dart implementation of the reviewed RIII
formulas.  Its baseline source and blob hashes
are pinned so that repository-side review can detect drift.

\subsection{Interpretation}

The validation supports four modest conclusions.
\begin{enumerate}
\item The generic implementation agrees with the theorem on a substantial
      finite sample.
\item The false connected-prefix conjecture is concretely guarded against.
\item The accumulated-support algorithm is genuinely different from the legacy
      last-touch algorithm.
\item Production correctness still depends on repository-native replay,
      restoration, invariant, and regression tests.
\end{enumerate}

\section{Consequences, comparisons, and limitations}
\label{sec:consequences}

\subsection{A new structural target}

Earlier conditional routes in the project parameterize exact computation by
frontier width, twist-run count, extremal Khovanov windows, transfer rank, or
hierarchy data.  The present result adds a parameter on the diagram-rewriting
side:
\[
  (\text{birth number},\ \text{unlock depth}).
\]
It is independent of homological width.  A diagram may have a difficult exact
scan but a small causal RIII certificate, or conversely a simple exact invariant
calculation but no short monotone simplification.

This separation suggests an adaptive policy: estimate local RIII branch growth
and exact-scan frontier growth, then spend a bounded amount on whichever has the
more favorable certified parameter.  Such a policy must preserve the existing
rule that abandoned work cannot affect verdicts.

\subsection{Why the result is stronger than a depth-only bound}

A naive global depth-$k$ RIII search has $O(N)$ choices at each level and hence
$N^{O(k)}$ behavior.  At $k=O(\log N)$ this is already
$N^{O(\log N)}$, quasi-polynomial, but with a large exponent and no recognition
of local causal structure.  More importantly, the current implementation does
not perform that complete global search because it would be too expensive.

The causal theorem decomposes the exponent.  Only births pay a factor of $N$;
all connected continuations pay $O(k)$.  When $b$ is constant and
$k=O(\log N)$, the exponent improves from $O(\log N)$ to
$O(\log\log N)$.  This is the main asymptotic gain.

\subsection{Why polynomial Reidemeister length is not enough}

Lackenby's polynomial move bound permits crossing increases and controls a much
broader class of sequences.  A polynomial total length $N^{11}$ inserted into
our enumeration formula would be useless.  The present theorem needs two finer
facts:
\begin{enumerate}
\item the neutral run between decreases is short;
\item only a few support-disjoint branches are born during that run.
\end{enumerate}
Proving these facts for all unknots would be a major additional theorem and may
be false for the strictly crossing-nonincreasing move language.

\subsection{Potential relation to planar separators}

The active RIII footprint contains $O(k)$ darts and induces a small subgraph of
the projection.  If this active region admits balanced separators that are
stable under local moves, one could attempt a divide-and-conquer search in the
style of frontier algorithms.  Such a theorem might replace the $N^b$ birth
factor by a smaller separator-dependent term.  At present, however, a birth can
occur anywhere, and no safe separator theorem for the dynamic RIII causal graph
is proved here.

\subsection{Potential relation to hierarchy certificates}

A short RIII causal trace is a two-dimensional diagram certificate; an essential
hierarchy is a three-dimensional complement certificate.  There may be a
transfer principle: a uniformly local hierarchy simplification could constrain
where diagrammatic births occur, while a small causal footprint might identify
a bounded part of a handle structure that needs modification.  The 2026
hierarchy algorithm makes this a timely direction, but no such bridge is
claimed.

\subsection{Lower-bound questions}

The birth parameter also supports negative results.  To prove that this route
cannot solve the unrestricted problem, one could construct unknot diagrams for
which every crossing-nonincreasing first-unlocking trace has either large depth
or large birth number.  The connectedness theorem says a shortest trace cannot
contain causally irrelevant branches, so a lower bound on births would be a
meaningful topological obstruction rather than a search artifact.

\subsection{Soundness boundary}

The preprocessor has one-sided semantics:
\[
  \text{successful replayable trace}\Rightarrow\Unk,
  \qquad
  \text{failure or timeout}\Rightarrow\Unknown.
\]
It cannot certify knottedness.  Knottedness remains the responsibility of sound
invariants or exact complement/homology methods.  This boundary should be
preserved in every CLI and serialized-certificate interface.

\subsection{Assessment of the research contribution}

Within its scope, the report delivers four concrete advances.
\begin{enumerate}
\item It replaces an empirical locality heuristic by an explicit two-parameter
      completeness theorem.
\item It proves a restricted quasi-polynomial and polynomial classification.
\item It identifies and repairs a subtle but natural false proof strategy.
\item It supplies code and a low-risk opt-in integration path.
\end{enumerate}
The unresolved step is geometric: establish useful upper bounds on births and
depth for broad, natural families of unknot diagrams, or design a different
certificate when those parameters are large.

\section{Further research questions and proposed program}
\label{sec:questions}

The following questions are ordered roughly from immediate engineering tasks to
longer-term topology.  They are phrased so that a negative answer would also be
informative.

\subsection{Immediate repository work}

\begin{enumerate}[label=\textbf{Q\arabic*.}]
\item \textbf{Find a genuine accumulated-support separator.}
      Does there exist a PD-code unknot for which the reviewed last-touch search
      fails at a given depth but the one-birth accumulated-footprint search
      succeeds at that depth?  An exhaustive search over small diagrams would
      turn Proposition~\ref{prop:last-vs-accumulated} from an abstract warning
      into a knot-specific regression fixture.

\item \textbf{Find a genuine two-birth RIII unlock.}
      Is there an unknot diagram whose shortest first-unlocking RIII trace has
      birth number at least two with the conservative dart footprint?  If so,
      minimize crossings and depth and add the result as an adversarial test.

\item \textbf{Measure the empirical parameter distribution.}
      On the existing random closures, scrambled unknots, hard survivors, and
      named examples, record the minimum successful pair $(b,k)$ under a fixed
      trial cap.  Separate successful and failed search costs, and publish raw
      traces rather than only aggregate counts.

\item \textbf{Incremental candidate maintenance.}
      Can legal RIII triangles meeting the active footprint be updated in
      amortized constant time per move, rather than rediscovered from nearby
      darts?  The dart involution changes only a bounded footprint, so a local
      invalidation table should be possible.

\item \textbf{Canonical commutation reduction.}
      Beyond sorting births, can independent connected moves be quotiented by a
      trace-monoid normal form?  A safe partial-order reduction could remove many
      permutations without changing completeness.
\end{enumerate}

\subsection{Parameterized topology}

\begin{enumerate}[label=\textbf{Q\arabic*.},resume]
\item \textbf{Bounded births for natural diagram classes.}
      Are reduced alternating unknot diagrams (necessarily trivial after
      reduction), positive braid-derived unknots, arborescent unknots, or
      diagrams of bounded projection-treewidth guaranteed to have bounded birth
      number during monotone simplification?

\item \textbf{Logarithmic depth from bounded width.}
      Does bounded treewidth, branchwidth, carving width, or pathwidth of the
      projection graph imply $k=O(\log N)$ after an appropriate choice of
      crossing-decreasing reduction order?  A balanced-separator proof would
      directly combine with the present theorem.

\item \textbf{Birth number versus dependency-tree width.}
      The support-overlap graph of a shortest trace is connected.  Relate the
      minimum birth number to standard directed or undirected parameters of this
      graph: number of sources under chronological orientation, pathwidth,
      vertex separation, or search number.  Which parameter admits the best
      deterministic enumeration?

\item \textbf{Lower bounds.}
      Construct families of unknot diagrams requiring
      $\Omega(\log N)$, $\Omega(N^\epsilon)$, or $\Omega(N)$ births or RIII
      depth under crossing-nonincreasing RI/RII/RIII simplification.  Such a
      family would sharply delimit this route.

\item \textbf{Reduction-order optimization.}
      The current simplifier takes available RI/RII moves greedily.  Can a
      different choice of decreasing move reduce the later birth/depth
      parameters?  Determine the complexity of finding an optimal reduction
      order and whether local confluence or matroid-like exchange principles
      apply.
\end{enumerate}

\subsection{Algorithmic improvements to the birth exponent}

\begin{enumerate}[label=\textbf{Q\arabic*.},resume]
\item \textbf{Color-coding independent births.}
      Can randomized or derandomized color-coding find $b$ disjoint causal birth
      footprints in $2^{O(b)}\poly(N)$ rather than $N^b$ time, while still
      accounting for their later joins?

\item \textbf{Terminal-rooted reverse search.}
      RIII is an involution.  Can one enumerate causal predecessors backward
      from potential RI/RII footprints, meeting a forward search in the middle?
      The obstacle is that the final diagram is unknown; a bounded local-state
      interface might overcome it.

\item \textbf{Representative families.}
      When several birth sets induce the same boundary behavior on an active
      footprint, can all but a small representative family be discarded?  This
      would parallel rank-based or connectivity-based dynamic programming.

\item \textbf{Planar localization.}
      Prove that a shortest birth footprint can be chosen inside a bounded-radius
      neighborhood of the eventual connecting action, or find a counterexample.
      A positive theorem would replace global $N$ choices by a function of $k$.

\item \textbf{Learning only as a scheduler.}
      Can a learned model predict which birth candidates to try first while the
      complete bounded search remains available underneath?  The learned layer
      must affect order only, never pruning or verdicts.
\end{enumerate}

\subsection{Combination with exact invariants}

\begin{enumerate}[label=\textbf{Q\arabic*.},resume]
\item \textbf{Khovanov-aware RIII ordering.}
      Can inexpensive frontier estimates predict which legal RIII move will
      reduce later Khovanov scan width, not merely crossing count?  The synthesis
      already observes that fewer crossings can increase scan work.

\item \textbf{Joint parameter theorem.}
      Suppose RIII birth/depth are moderately large but the resulting diagram
      has a small exact-scan frontier.  Derive a combined bound that branches on
      RIII only until a certified frontier threshold is reached, then switches
      to exact contraction.

\item \textbf{Twist-run compression after causal simplification.}
      Does a small RIII causal footprint preserve or reveal a low signed-braid
      run count?  If so, the twist-compression theorem and the present theorem
      could be composed rather than run as unrelated alternatives.

\item \textbf{Hierarchy transfer.}
      Relate a bounded RIII causal region to a bounded modification of a handle
      hierarchy or boundary pattern.  This could provide a three-dimensional
      certificate that replaces a large remaining diagram search.
\end{enumerate}

\subsection{Formal verification}

\begin{enumerate}[label=\textbf{Q\arabic*.},resume]
\item \textbf{Lean formalization of the abstract theorem.}
      Formalize strong disjoint commutation, causal connectedness, birth
      front-loading, and the counting bound.  The Boolean prototype supplies a
      finite executable semantics, but the proof should be independent of it.

\item \textbf{Verified dart footprint.}
      Prove that the concrete \code{apply\_r3} function writes only the claimed
      24-dart footprint and is an involution on every validated diagram.

\item \textbf{Proof-carrying traces.}
      Specify a stable JSON schema for RIII/RI/RII traces, source hashes, budgets,
      and replay results.  Make the replay kernel the only trusted component for
      positive simplification verdicts.
\end{enumerate}

\subsection{A focused next milestone}

The most informative next experiment is not to raise depth globally.  It is to
implement opt-in $b=1$ accumulated support, preserve the current budget, and run
paired tests on the exact 117-depth corpus recorded by the synthesis.  Three
outcomes are all useful:
\begin{itemize}
\item new successes establish that last-touch locality was a real completeness
      bottleneck;
\item no new successes but low overhead suggests the stronger search can safely
      replace the heuristic;
\item high failed-search overhead identifies the need for partial-order
      reduction or a better scheduler before increasing $b$.
\end{itemize}

\appendix
\section{Reproducibility, audit trail, and exact commands}
\label{sec:reproducibility}

\subsection{Reviewed source pins}

\begin{longtable}{@{}p{0.28\textwidth}p{0.66\textwidth}@{}}
\toprule
Item & Pin \\
\midrule
\endhead
ProveIt repository & \code{VladimirReshetnikov/ProveIt} \\
Reviewed commit & \hashcode{23893144e3aafd074a23cfb633222ab7e03e6d2e} \\
RIII source path & \code{Topology/UnknotRecognition/fast/fastunknot/simplify.py} \\
RIII source blob & \hashcode{2fcaabb51eaa45cf84bdee32cd1b4e18acf2c478} \\
OpenAI math README blob & \hashcode{8ac0cb7dbf1f38ad3b7480310d4e1d9482f0cfe8} \\
Date of review & October 7, 2026, America/Los\_Angeles \\
\bottomrule
\end{longtable}

\subsection{Bundle commands}

From the artifact root:

\begin{lstlisting}[language=bash]
python -m unittest discover -s tests -v
python experiments/exhaustive_validation.py \
  --systems 5000 \
  --sites 6 \
  --moves 9 \
  --reductions 3 \
  --depth 6 \
  --output results/validation.json
python experiments/exhaustive_catalog.py \
  --depth 4 \
  --moves-per-system 3 \
  --output results/exhaustive_catalog.json

cd article
./build.sh
\end{lstlisting}

The JSON output is deterministic because each generated system uses its integer
index as the random seed.

\subsection{Claims ledger}

The machine-readable file \code{claims/claims\_ledger.yaml} distinguishes:
\begin{itemize}
\item proved manuscript claims;
\item computational checks and their scope;
\item the unexecuted production integration;
\item the explicitly unproved general quasi-polynomial goal;
\item the absence of an exhaustive literature-priority claim.
\end{itemize}
The file \code{claims/research\_log.md} records the rejected connected-prefix
conjecture, the state-history and inverse-pruning implementation corrections,
and the deterministic-policy boundary of the recognition theorem.

\subsection{How to audit the key proof mechanically}

A lightweight independent audit can proceed as follows.
\begin{enumerate}
\item Check that every neutral rule's guard reads only its declared support and
      that disjoint supports imply commuting bit toggles.
\item For every unrestricted BFS witness, compute its births, front-load them,
      replay, and truncate at the first terminal state.
\item Run the birth-front search with the original witness's birth count and
      depth.  It must find a witness no longer than the oracle witness.
\item Build the support-overlap graph of the oracle witness plus its terminal
      reduction.  If it is disconnected, delete a nonterminal component and
      obtain a shorter witness, contradicting BFS minimality.
\end{enumerate}
The validation script performs these checks on generated finite systems.

\subsection{PDF verification}

The compiled PDF should be rendered to images with the repository-independent
PDF tooling and inspected for clipped equations, code blocks, table overflow,
and broken glyphs.  The build in this bundle uses only standard TeX packages and
ASCII punctuation in source text where renderer compatibility can be fragile.

\subsection{Repository integration boundary}

The patch is a review artifact.  It should not be merged merely because the
abstract tests pass.  The repository-side maintainer should:
\begin{enumerate}
\item verify the baseline blob or rebase the patch by inspection;
\item add repository-native unit tests for state restoration and replay;
\item run every maintained test and pinned-source check;
\item benchmark failed searches under the same paired methodology used by the
      synthesis;
\item leave clustered mode opt-in until the evidence supports a default change.
\end{enumerate}

\begin{thebibliography}{99}

\bibitem{BarNatan2005}
D. Bar--Natan,
\newblock Khovanov's homology for tangles and cobordisms,
\newblock \emph{Geometry \& Topology} 9 (2005), 1443--1499.

\bibitem{Dynnikov2006}
I. A. Dynnikov,
\newblock Arc-presentations of links: monotonic simplification,
\newblock \emph{Fundamenta Mathematicae} 190 (2006), 29--76;
\newblock arXiv:math/0208153.

\bibitem{HassLagarias2001}
J. Hass and J. C. Lagarias,
\newblock The number of Reidemeister moves needed for unknotting,
\newblock \emph{Journal of the American Mathematical Society} 14 (2001),
399--428;
\newblock arXiv:math/9807012.

\bibitem{HenrichKauffman2011}
A. Henrich and L. H. Kauffman,
\newblock Unknotting unknots,
\newblock arXiv:1006.4176v4 (2011).

\bibitem{KronheimerMrowka2011}
P. B. Kronheimer and T. S. Mrowka,
\newblock Khovanov homology is an unknot-detector,
\newblock \emph{Publications Mathematiques de l'IHES} 113 (2011), 97--208.

\bibitem{Lackenby2015}
M. Lackenby,
\newblock A polynomial upper bound on Reidemeister moves,
\newblock \emph{Annals of Mathematics} 182 (2015), 491--564;
\newblock arXiv:1302.0180.

\bibitem{Lackenby2026}
M. Lackenby,
\newblock Incompressible surfaces, hierarchies and unknot recognition,
\newblock arXiv:2607.23350v1 (2026).

\bibitem{OpenAIMath2026}
OpenAI,
\newblock Mathematical manuscripts and supporting proof artifacts,
\newblock \url{https://github.com/openai/math}, reviewed October 7, 2026.

\bibitem{ProveIt2026}
V. Reshetnikov,
\newblock ProveIt: UnknotRecognition,
\newblock \url{https://github.com/VladimirReshetnikov/ProveIt/tree/main/Topology/UnknotRecognition},
reviewed at commit
\code{23893144e3aafd074a23cfb633222ab7e03e6d2e}, October 7, 2026.

\end{thebibliography}


\end{document}
